\section*{Chapitre \ref{ch:iv1:timed-test-2} --- Épreuve chronométrée 2}
\markboth{\protect\omTocMark{}{Réponses~: Épreuve chronométrée 2}}{\protect\omTocMark{}{Réponses~: Épreuve chronométrée 2}}

\begin{solution}{iq:iv1:timed-test-2:1}
La moyenne est de $\frac74$ boules.

Utiliser des variables indicatrices\index{variables indicatrices}. Une boule rouge n'est laissée de côté que si les 3 boules bleues sortent toutes avant elle. Parmi cette boule rouge et les 3 boules bleues, tous les ordres sont équiprobables, donc la probabilité que la rouge sorte en dernier est $\frac14$. De même, une boule bleue n'est laissée de côté que si les 5 rouges sortent toutes avant elle, avec une probabilité $\frac16$.

En sommant, $5 \times \frac14 + 3 \times \frac16 = \frac54 + \frac12 = \frac74$.

\iqwrong{Conditionner selon la couleur qui s'épuise la première, puis chercher le nombre de boules restantes dans chaque cas. On peut y parvenir, mais c'est beaucoup plus long, et les pondérations des cas sont faciles à se tromper.}
\end{solution}

\begin{solution}{iq:iv1:timed-test-2:2}
La moyenne est de $19.5$ lancers.

Un lancer donne 7 avec une probabilité de $\frac16$ et 11 avec une probabilité de $\frac{1}{18}$. Soient $A$ et $B$ les temps d'attente de chacun. Le temps nécessaire pour voir les deux est $\max(A, B) = A + B - \min(A, B)$\index{inclusion-exclusion}.

$\mathbb{E}A = 6$ et $\mathbb{E}B = 18$. Le minimum est le temps d'attente de l'un ou l'autre des deux totaux, qui a une probabilité $\frac16 + \frac{1}{18} = \frac29$ par lancer, donc $\mathbb{E}\min = 4.5$. La réponse est $6 + 18 - 4.5 = 19.5$.

Notons que le résultat n'excède que de $1.5$ l'attente de l'événement le plus rare pris seul, car l'événement fréquent survient en général en chemin.

\iqwrong{Additionner 6 et 18. Les deux attentes se chevauchent, car un même lancer compte pour les deux à la fois.}
\end{solution}

\begin{solution}{iq:iv1:timed-test-2:3}
Le prix est $\frac{20}{3} \approx 6.67$, et la couverture consiste en $\frac23$ d'action.

L'option d'achat paie 20 ou 0. Une position de $\Delta$ actions varie de $30\Delta$ entre les deux états, donc pour égaler l'écart de 20, $\Delta = \frac23$\index{réplication}.

Avec $\frac23$ d'action, le portefeuille vaut 80 ou 60. Emprunter 60 à taux nul donne exactement 20 ou 0, donc l'option d'achat coûte $\frac23 \times 100 - 60 = 6.67$.

Une seconde méthode~: la probabilité risque-neutre de l'état haussier résout $120q + 90(1 - q) = 100$, donc $q = \frac13$, et le prix vaut $\frac13 \times 20$. La probabilité réelle de l'état haussier n'intervient nulle part.

\iqfollow{Que feriez-vous si l'option d'achat était cotée à 7~?}{On la vend, on achète $\frac23$ d'action et on emprunte le reste. La position rapporte alors 33 centimes dans les deux états, sans aucun risque.}
\end{solution}

\enlargethispage{-3\baselineskip}%
\begin{solution}{iq:iv1:timed-test-2:4}
La probabilité est $\frac16$.

La frontière est $\sqrt{y} = 1 - \sqrt{x}$, c'est-à-dire $y = (1 - \sqrt{x})^2$, donc l'événement est la région située sous cette courbe dans le carré unité. Son aire vaut
\[
\int_0^1(1 - \sqrt{x})^2\,dx = \int_0^1(1 - 2\sqrt{x} + x)\,dx = 1 - \frac43 + \frac12 = \frac16.
\]
Une seconde méthode\index{changement de variables}~: on pose $U = \sqrt{X}$ et $V = \sqrt{Y}$. Comme $P(U \le u) = u^2$, chacune a pour densité $2u$ sur $[0, 1]$, et la réponse est $\iint_{u + v < 1}4uv\,du\,dv = 4 \cdot \frac{1}{24} = \frac16$.

\iqwrong{Répondre $\frac12$, en prenant la frontière pour la droite $x + y = 1$. Sur $[0, 1]$, on a $\sqrt{x} \ge x$, donc l'événement impose $X + Y < 1$ et ne couvre qu'une partie de ce triangle~: la courbe s'infléchit vers l'origine.}
\end{solution}

\begin{solution}{iq:iv1:timed-test-2:5}
La probabilité est $\binom{10}{5}^2/\binom{20}{10} = 63\,504/184\,756 \approx 0.344$.

Toute suite particulière de 20 lancers de pièce équilibrée a la même probabilité, donc, sachant qu'il y a 10 piles au total, chacune des $\binom{20}{10}$ façons de placer ces piles est équiprobable. Obtenir exactement 5 piles parmi les 10 premiers lancers revient à choisir 5 places parmi les 10 premières et 5 parmi les 10 dernières, ce qui donne $\binom{10}{5}^2$ façons\index{loi hypergéométrique}.

Numériquement, $\binom{2n}{n} \approx 4^n/\sqrt{\pi n}$ ramène le rapport à $\sqrt{10\pi}/(5\pi) \approx 0.357$, proche de la valeur exacte $0.344$.

\iqwrong{Répondre $\binom{10}{5}/2^{10} \approx 0.246$, la probabilité d'obtenir 5 piles en 10 lancers sans condition. Fixer le total lie les deux moitiés~: plus de piles au début imposent moins à la fin, ce qui resserre la loi du nombre de piles de la première moitié.}

\iqfollow{Comment comparez-vous ces deux probabilités en général, pour de grands échantillons~?}{La probabilité conditionnelle est plus grande d'un facteur d'environ $\sqrt{2}$, puisque la variance du nombre de piles dans la première moitié est réduite de moitié. Ici, $0.344/0.246 = 1.40$.}
\end{solution}

\begin{solution}{iq:iv1:timed-test-2:6}
L'espérance du nombre de lancers est $\frac83$.

Soit $N$ le nombre de lancers, donc $P(N = k) = 2^{-k}$ et $P(N = 2j) = 4^{-j}$. Leur somme vaut $\frac13$, la probabilité d'un nombre pair de lancers. Sachant qu'il est pair, $P(N = 2j) = 3 \cdot 4^{-j}$, donc $N/2$ suit une loi géométrique de paramètre $\frac34$ et d'espérance $\frac43$\index{conditionnement sur un événement}. Ainsi $N$ a pour espérance conditionnelle $\frac83$.

Un calcul direct confirme~: $\sum_j 2j\,4^{-j} = \frac89$, et $\frac89 \big/ \frac13 = \frac83$.

\iqwrong{Répondre 4, en doublant l'espérance inconditionnelle de 2, puisque $N$ vaut le double de $N/2$. Sachant qu'il est pair, $N/2$ n'a pas la même loi que $N$~: la valeur 1 a pour probabilité $\frac34$, et non $\frac12$.}

\iqfollow{Quel résultat obtenez-vous pour une pièce dont la probabilité de pile est $p$~?}{$2/(1 - (1 - p)^2)$. Il tend vers 2 lorsque $p$ tend vers 1, comme il se doit, puisque 2 est le plus petit nombre pair de lancers possible.}
\end{solution}

\begin{solution}{iq:iv1:timed-test-2:7}
La probabilité est $\frac{12}{38} = \frac{6}{19} \approx 0.316$.

Les 39 cartes que l'on ne voit pas forment trois mains de 13. Placer les deux as manquants un à un parmi ces 39 emplacements\index{placement séquentiel}. Le premier peut aller n'importe où. Le second doit prendre l'un des 12 autres emplacements de cette main, parmi les 38 restants.

Un comptage des paires donne la même valeur~: sur les $\binom{39}{2} = 741$ paires d'emplacements, $3\binom{13}{2} = 234$ sont dans une même main, et $\frac{234}{741} = \frac{6}{19}$.

\iqwrong{Répondre $\frac13$, puisqu'il y a trois mains où placer le second as. Une main contenant déjà un as n'a que 12 emplacements libres, contre 13 dans chacune des autres. Avec des mains de $h$ cartes, la probabilité vaut $(h - 1)/(3h - 1)$~: elle est nulle pour $h = 1$, et tend vers $\frac13$ seulement pour de grandes valeurs de $h$.}
\end{solution}

\begin{solution}{iq:iv1:timed-test-2:8}
La moyenne est de $\frac{20}{19} \approx 1.05$ couples.

Utiliser une variable indicatrice\index{variables indicatrices} pour chaque couple. Fixer la place d'un des conjoints. L'autre conjoint a la même probabilité d'occuper l'un quelconque des 19 sièges restants, et exactement 2 d'entre eux sont voisins. Chaque couple est donc assis côte à côte avec la probabilité $\frac{2}{19}$.

Par linéarité, la moyenne vaut $10 \times \frac{2}{19} = \frac{20}{19}$. Les variables indicatrices sont dépendantes, mais la linéarité ne s'en soucie pas.

La moyenne vaut $2n/(2n - 1)$ pour $n \ge 2$ couples~: elle tend vers 1 lorsque $n$ croît, et le nombre de couples réunis tend vers une loi de Poisson de moyenne 1.

\iqwrong{Chercher la loi du nombre de couples voisins. Cela demande l'inclusion--exclusion, alors que la moyenne ne demande qu'une ligne.}
\end{solution}

\begin{solution}{iq:iv1:timed-test-2:9}
$\rho = 0.24/0.99 \approx 0.242$.

La variance d'un panier de $N$ titres à pondération égale est\index{corrélation implicite}
\[
\sigma_I^2 = \sigma^2\left(\frac1N + \left(1 - \frac1N\right)\rho\right).
\]
Avec $\sigma_I^2/\sigma^2 = (15/30)^2 = 0.25$ et $N = 100$~: $0.25 = 0.01 + 0.99\rho$, donc $\rho = 0.24/0.99$.

Sans la correction pour $N$ fini, la réponse rapide est $0.25$, ce qui est déjà proche. Pour $N$ grand, le rapport des deux variances se réduit à la corrélation, et c'est ainsi qu'un desk lit une corrélation moyenne à partir d'un indice et de ses composantes.

Contrôle de cohérence~: pour $\rho = 1$, la volatilité de l'indice est de 30~\%, et pour $\rho = 0$, de 3~\%.

\iqfollow{Les options sur l'indice impliquent 15~\%, et celles sur les actions 30~\%. Un trader pense que la corrélation réalisée sera de $0.15$. Quelle volatilité en déduisez-vous pour l'indice~?}{$30~\% \times \sqrt{0.01 + 0.99 \times 0.15} = 11.9~\%$, inférieure aux 15~\% implicites. C'est la stratégie que les desks de dispersion envisagent.}
\end{solution}

\begin{solution}{iq:iv1:timed-test-2:10}
La probabilité vaut $\frac{5 - \pi}{4} \approx 0.4646$.

Le rapport s'arrondit à 0 lorsque $X/Y < \frac12$, avec une probabilité $\frac14$ (un triangle de base $\frac12$ et de hauteur 1).

Pour $r \ge 1$, $P(X/Y > r) = 1/(2r)$. Ainsi le rapport s'arrondit à $2k$, pour $k \ge 1$, avec une probabilité
\[
\frac12\left(\frac{1}{2k - \frac12} - \frac{1}{2k + \frac12}\right) = \frac{1}{4k - 1} - \frac{1}{4k + 1}.
\]
La somme sur $k$ vaut $\frac13 - \frac15 + \frac17 - \frac19 + \dots$, qui, par la série de Leibniz\index{série de Leibniz}, vaut $1 - \frac{\pi}{4}$.

En ajoutant $\frac14$, on obtient $\frac54 - \frac{\pi}{4}$. Une simulation donne $0.465$. Elle est inférieure à $\frac12$, car la valeur impaire 1 seule porte environ $0.42$ de la masse, davantage que la valeur paire 0.

\iqwrong{Utiliser $P(X/Y > r) = 1/r$. La région $X > rY$ dans le carré unité est un triangle de base 1 et de hauteur $1/r$, donc il y a un facteur $\frac12$.}
\end{solution}

\begin{solution}{iq:iv1:timed-test-2:11}
La probabilité vaut environ $0.736$.

On découpe selon le score auquel la partie se termine. Avec $p = 0.6$ et $q = 0.4$~:

Gagner à zéro, à 15 et à 30~: $p^4(1 + 4q + 10q^2) = 0.1296 \times 4.2 = 0.544$.

Atteindre l'égalité, 3--3~: $\binom63p^3q^3 = 20 \times 0.216 \times 0.064 = 0.276$.

Depuis l'égalité, deux points décident ou ramènent à l'égalité. Gagner deux points avec $p^2$, en perdre deux avec $q^2$, donc la probabilité de gagner depuis l'égalité est\index{série géométrique}
\[
\frac{p^2}{p^2 + q^2} = \frac{0.36}{0.52} = 0.692.
\]
Total~: $0.544 + 0.276 \times 0.692 = 0.736$. Une récurrence exacte sur tous les scores confirme ce résultat.

Un edge de $0.6$ par point devient un edge de $0.74$ par partie, et c'est ainsi que de petits edges se cumulent au fil des parties répétées.

\iqwrong{Oublier les coefficients, ou écrire $\binom{6}{3}$ à tort pour zéro, 15 et 30. Le dernier point doit être gagné, donc les comptes sont $\binom{3}{0}$, $\binom{4}{1}$ et $\binom{5}{2}$.}
\end{solution}

\begin{solution}{iq:iv1:timed-test-2:12}
La somme vaut 6.

Soit $N$ le nombre de lancers jusqu'au premier pile, de sorte que $P(N = n) = 2^{-n}$. La somme vaut $\mathbb{E}N^2$\index{loi géométrique}. Pour une variable géométrique avec $p = \frac12$, l'espérance vaut 2, et la variance vaut $(1 - p)/p^2 = 2$. Donc $\mathbb{E}N^2 = 2 + 4 = 6$.

Une voie directe~: à partir de $\sum x^n = 1/(1 - x)$, en dérivant puis en multipliant deux fois par $x$, on obtient $\sum n^2x^n = x(1 + x)/(1 - x)^3$. En $x = \frac12$, cela vaut $\frac12 \cdot \frac32 \cdot 8 = 6$. Une somme partielle le confirme.

\iqfollow{Que vaut $\sum n^3/2^n$~?}{26. En dérivant une fois de plus, on obtient $x(1 + 4x + x^2)/(1 - x)^4 = \frac12 \cdot \frac{13}{4} \cdot 16 = 26$.}

\iqwhy{Ramener une série à un moment fait gagner du calcul et donne une seconde façon de vérifier la réponse.}
\end{solution}
